\section{Mutation Operators}
\begin{longtable}{p{.27\linewidth}p{.63\linewidth}}
\toprule
Operator & Required witness \\
\midrule
Cardinality sum to product & Unequal choice with three nonunit branches \\
Omit child count & Last branch contributes positive cardinality \\
Prefix off by one & Rank exactly at a branch boundary \\
Reverse branch order & Rank zero selects the declared first branch \\
Exclusive range stop & Upper endpoint must rank successfully \\
Shift lower bound & Lower endpoint must produce zero local offset \\
Reset rank accumulator & Later branch must retain its prefix \\
Multiply rank offset & Range step uses child block exactly once \\
Unrank uses non-strict boundary & Boundary belongs to the following branch \\
Wrong boundary child & Second transition selects third branch \\
Boolean accepted as integer & Typed integer field rejects Boolean \\
Unicode normalization removed & Composed and decomposed strings hash identically \\
Hash ignores branch order & Ordered branch reversal changes identity \\
Overflow wraps & Cardinality beyond native bound is rejected \\
Missing target accepted & Graph validation rejects unresolved edge \\
Duplicate label accepted & Canonically duplicate branch is rejected \\
Trailing token ignored & Terminal requires complete token consumption \\
Unreachable node accepted & Strict source schema rejects unreachable content \\
\bottomrule
\end{longtable}
Python applies all operators. Native C and Rust apply the cardinality, range, rank, boundary, and overflow subset because they consume already validated canonical IR. The raw reports preserve applicability and compile status.
